\chapter*{\centering ВВЕДЕНИЕ}
\addcontentsline{toc}{chapter}{ВВЕДЕНИЕ}




\Edit{ВВЕДЕНИЕ
должно
содержать
оценку
современного
состояния
решаемой научно-исследовательской или научно-технической проблемы,
основание и исходные данные для разработки темы; обоснование выбора темы,
определяемой ее актуальностью; формулировка проблемы и круга вопросов,
необходимых
для
ее
решения;
цель
работы
с
ее
разделением
на
взаимосвязанный комплекс задач, подлежащих решению, для раскрытия темы;
указываются объект исследования или разработки, определяются методы
исследования, дается краткий обзор базы исследования и используемых
источников.}

В задачах городской логистики может возникнуть необходимость поиска маршрута по нескольким критериям. Например, при мультимодальных поездках оптимизируют время с учётом расписаний и пересадок~\cite{teslya2016multimodal}, а при проектировании транспортной сети сопоставляют стоимость, время и надёжность инфраструктуры~\cite{medvedeva2024rail}. 


Целью данной работы является разработка метода поиска компромиссного пути в графе городской логистики на основе генетического алгоритма.

Для достижения поставленной цели необходимо выполнить перечисленные ниже задачи.

\begin{enumerate}
    \item Провести анализ предметной области мультикритериального поиска путей. Рассмотреть известные генетические алгоритмы. Провести сравнительный анализ рассмотренных алгоритмов. Формализовать постановку задачи в виде функциональной модели.
    \item Разработать метод поиска компромиссного пути в графе городской логистики на основе генетического алгоритма. Описать основные особенности разрабатываемого метода. Сформулировать и описать ключевые шаги метода в виде схем алгоритмов. Описать структуры данных, используемые в алгоритмах.
    \item Осуществить выбор программных средств реализации предложенного метода. Разработать программное обеспечение, реализующее описанный метод. Описать взаимодействие пользователя с программным обеспечением.
    \item Провести исследование зависимости времени поиска компромиссного пути от количества вершин в графе, критериев поиска.
\end{enumerate}
